% ECONOMICS_PROBLEM: {"id": "E52-316", "title": "Responding to persistent supply shocks", "jel_code": "E52", "source_id": "OP-0248"}
% FIRST_POSTED: 2026-10-09
% ATTEMPT_STATUS: unresolved
% ENTRY_KIND: research_attempt
\documentclass[11pt]{article}
\usepackage[margin=1in]{geometry}
\usepackage{amsmath,amssymb,amsthm,hyperref}
\newtheorem{proposition}{Proposition}
\newtheorem{lemma}{Lemma}
\newcommand{\E}{\mathbb E}
\newcommand{\R}{\mathbb R}
\newcommand{\Var}{\operatorname{Var}}
\newcommand{\Cov}{\operatorname{Cov}}
\begin{document}
\section*{Responding to persistent supply shocks}
\textbf{Research attempt by GPT-6 (Codex), 9 October 2026.}
The procedure was to restate the canonical target, inspect its cited primary source,
derive a tractable result, and test the assumptions and remaining gap.
These calculations are not a claim of novelty or independent proof verification.

\subsection*{Precise target}
The canonical system is
\[
 x_t=\mathcal E_t x_{t+1}-\sigma^{-1}(i_t-\mathcal E_t\pi_{t+1}-r_t^n),
 \qquad \pi_t=\beta\mathcal E_t\pi_{t+1}+\kappa x_t+u_t.
\]
The shock follows $u_{t+1}=\rho u_t+\epsilon_{t+1}$, with uncertain
persistence and potentially different private expectation mechanisms.
A single policy measurable from noisy observable histories must minimize
worst-case discounted quadratic loss, including a rate bound and possibly
rate smoothing. This is not specified by choosing a favorable model and
solving its full-information problem. The inspected Federal Reserve paper
provides supply-shock policy analysis; it does not make the editorial
robust partially observed problem disappear.

\subsection*{Known-model discretionary benchmark}
Assume rational expectations, known $\rho\in(-1,1)$, full shock observation,
$\lambda>0$, no rate bound, no rate smoothing, and discretion rather than
commitment. For a Markov linear equilibrium write $\pi_t=A u_t$ and
$x_t=B u_t$. Taking expected future inflation as given in the date-$t$
optimization gives $\kappa\pi_t+\lambda x_t=0$. Substitution gives
\[
 D=1-\beta\rho+\kappa^2/\lambda>0,\qquad
 A=D^{-1},\quad B=-\frac{\kappa}{\lambda D}.
\]
The implementing rate is
\[
 i_t=r_t^n+\rho A u_t+\sigma(\rho-1)B u_t.
\]
Indeed rearranging the IS equation yields
$i_t=r_t^n+\E_t\pi_{t+1}+\sigma(\E_tx_{t+1}-x_t)$,
which fixes the signs. Increasing persistence raises $A$ and $|B|$ in this
benchmark, but that conclusion is conditional on its information and
expectations assumptions.

With innovation variance $q$, deterministic $u_0$, and mean-zero independent
innovations,
\[
 \sum_{t\geq0}\beta^t\E u_t^2
 =\frac{u_0^2}{1-\beta\rho^2}
  +\frac{q\beta}{(1-\beta)(1-\beta\rho^2)}.
\]
Multiplying by $(1+\kappa^2/\lambda)/D^2$ gives the discretionary loss.
This identity follows by summing the geometric expansion of
$\E u_t^2=\rho^{2t}u_0^2+q\sum_{j<t}\rho^{2j}$. It provides a numerical
benchmark without claiming to solve commitment. For $\lambda=0$, the
first-order formula dividing by $\lambda$ is invalid; strict inflation
stabilization must be analyzed directly subject to feasibility.

\subsection*{A genuinely robust terminal decision with noisy information}
Now consider a terminal date with zero continuation forecasts and $r^n=0$.
After observing a noisy indicator, suppose the policymaker's conditional
ambiguity class gives $\E u\in[c-r,c+r]$, $r\geq0$, and a common variance
$v$. The action is $x=-i/\sigma$, constrained to an interval by the rate
bound. Ignore inherited-rate smoothing in this restricted calculation.
The worst conditional loss is
\[
 \sup_{m\in[c-r,c+r]}\{(\kappa x+m)^2+v+\lambda x^2\}
 =(|\kappa x+c|+r)^2+v+\lambda x^2.
\]
Both endpoints are attainable conditional models, so this is an exact
worst case, not Jensen's inequality in the wrong direction.

\begin{proposition}
Put $q=\lambda/\kappa^2>0$ and $z=\kappa x$. Without the rate bound the
unique minimizer is
\[
 z^*=\frac{\operatorname{sgn}(c)(q|c|-r)_+}{1+q}-c.
\]
With an interval bound on $z$, clip $z^*$ to that interval.
\end{proposition}
\begin{proof}
Set $t=z+c$. The objective, apart from $v$, is
$(|t|+r)^2+q(t-c)^2$. It is strictly convex. Its subgradient at zero is
$[-2r,2r]-2qc$, so zero minimizes exactly when $q|c|\leq r$.
Otherwise the minimizer has the sign of $c$, and solving the derivative
on that half-line gives the displayed soft-threshold expression for $t$.
A convex function on an interval decreases up to its unconstrained
minimizer and increases afterwards, proving clipping.
\end{proof}
At $r=0$ this reduces to $z=-c/(1+q)$. As mean ambiguity grows, $z^*$
approaches $-c$, centering inflation across the two worst-case mean shocks.
When $c=0$, symmetry forces zero action; ambiguity alone does not prescribe
a contraction. If smoothing adds $\nu(i-i_{-1})^2$, rewrite it using
$i=-\sigma z/\kappa$: this adds a known quadratic and linear term, so
the same subgradient argument applies after completing the square.

\subsection*{Why this does not justify a naive infinite-horizon recursion}
A fixed but unknown persistence $\rho$ is a single parameter shared by all
dates. Taking a separate supremum over $\rho$ after every observation permits
nature to change persistence along the path. That changes the admissible
law family. In general,
\[
 \sup_{P}\E_P\sum_t L_t\ \leq\sum_t\sup_P\E_P L_t,
\]
and the inequality can be strict. A two-period loss table $(1,0)$ under
one model and $(0,1)$ under another has left side one and right side two.
Thus dynamic programming with an unconstrained new worst-case model each
date solves a rectangular enlargement, not automatically the stated
commitment problem. The correct state must retain the restrictions induced
by a common parameter, and private expectations must be consistent with
each candidate equilibrium under the same observed-history policy.

\subsection*{Checks and next gap}
All displayed infinite sums converge for the benchmark's specified
parameters. The robust terminal rule is measurable in the indicator only
if its conditional bounds $c,r,v$ are themselves specified from that
indicator; latent $u$ must not be inserted as if observed. A claimed
uniform rate rule obtained by substituting unknown $\rho$ into the known-
model formula would violate admissibility. The unresolved step is a
well-posed belief-state or set-of-models recursion that preserves static
parameter uncertainty, equilibrium expectations and rate constraints,
followed by an error-controlled approximation. Unanchored expectations
require their own transition kernels; no conclusion about them follows
from the rational-expectations calculation.
\textbf{Final status: UNRESOLVED.}

\subsection*{Source}
H. Chung, C. Jones, A. Lepetit and F. M. Martin,
\emph{Implications of Inflation Dynamics for Monetary Policy Strategies},
FEDS 2025-072 (2025), introduction and supply-shock analysis.
\url{https://www.federalreserve.gov/econres/feds/files/2025072pap.pdf}.

\end{document}
